% Original Prism Web Server Issue skeleton for Nucleic Acids Research.
% Uses the unmodified OUP 1.5 class and its Modern Large page design.
% The sections below are writing suggestions, not an official fixed outline.
\documentclass[unnumsec,webpdf,modern,large]{oup-authoring-template}
\usepackage{booktabs,tabularx,siunitx}
\newcommand{\societylogo}{}
\sisetup{uncertainty-mode=separate}

\begin{document}
\journaltitle{Nucleic Acids Research}
\DOI{}
\copyrightyear{}
\pubyear{}
\access{Draft manuscript}
\appnotes{Web Server Issue}
\firstpage{1}

\title[Tool name and short title]{Tool name: concise description of the web application}
\author[1,$\ast$]{First Author}
\author[1,2]{Second Author}
\address[1]{\orgdiv{Department}, \orgname{First institution},
  \orgaddress{\state{City}, \postcode{Postal code}, \country{Country}}}
\address[2]{\orgdiv{Department}, \orgname{Second institution},
  \orgaddress{\state{City}, \postcode{Postal code}, \country{Country}}}
\corresp[$\ast$]{Corresponding author.
  \href{mailto:first.author@example.org}{first.author@example.org}}

% One English paragraph, at most 200 words; include the actual application URL.
% example.org is reserved demonstration content, not a functioning web server.
\abstract{Replace this paragraph with the biological problem, intended
users, input, computation, output and verified value of the application.
Summarize the evaluation and the distinction from comparable tools.
State the actual application URL in the abstract. The reserved example
\url{https://example.org/tool} must be replaced; it does not identify a
working server. This starter reports no implemented service, validation
experiment or scientific finding.}
\maketitle
\thispagestyle{plain}
\pagestyle{plain}

\section{Introduction}\label{sec:introduction}
Explain the biological need and the intended audience. Identify related
tools and the specific benefit of the proposed application. For an update,
describe what has substantively changed since the previous publication.
The public reference below tests bibliography mechanics only; retain it
only if relevant to the work~\cite{wilkinson2016fair}.

\section{Materials and Methods}\label{sec:methods}
\subsection{Algorithms and implementation}
Describe the methods, software dependencies and versions, implementation
and computational workflow. Explain any new method in enough detail to
reproduce it; do not relegate essential procedures to the supplement.

\subsection{Validation design}
Describe the actual benchmark data, comparison tools and settings,
independent evaluation and statistical procedures. For a predictive tool,
separate validation data from training and model selection. Explain
applicable domain benchmarks and limitations rather than claiming
unperformed validation.

Equation~\eqref{eq:precision} is a formatting example, not a measured
performance result. Define its terms and relevance, or remove it.
\begin{equation}
 p=\frac{n_{\mathrm{correct}}}{n_{\mathrm{returned}}}.
 \label{eq:precision}
\end{equation}

\section{Input, Processing and Output}\label{sec:workflow}
Give the working application URL in the body as well as the abstract.
The reserved example \url{https://example.org/tool} must be replaced.
Describe accepted inputs, limits, parameters, processing, result formats
and interpretation. Explain the sample-data mechanism, help and tutorial
that users can actually access. Document how submitted data are kept
private and how users retrieve results, without promising unimplemented
safeguards.

\begin{figure}[t]
\centering
\fbox{\parbox[c][33mm][c]{0.87\linewidth}{\centering
Input to output workflow\\[5pt]Replace this synthetic slot}}
\caption{Synthetic workflow slot. Replace with an original figure that
explains a real computation and how the returned results are interpreted.}
\label{fig:workflow}
\par\smallskip\small\raggedright
\textbf{Alt text:} Replace with a concise description of the actual workflow.
\end{figure}

\section{Results and Use Cases}\label{sec:results}
Report verified evaluation results and comparison with appropriate
existing methods. Describe one or more use cases showing the biological
insight the application can provide. Figure~\ref{fig:workflow} is a layout
slot only; Table~\ref{tab:evaluation} contains no observed result.

\begin{table}[t]
\centering\small
\caption{Invented values illustrating an evaluation table.}
\label{tab:evaluation}
\begin{tabularx}{\linewidth}{@{}>{\raggedright\arraybackslash}X
 S[table-format=2.0] S[table-format=1.3(3)]@{}}
\toprule
Setting & {Count} & {Estimate $\pm$ SD}\\
\midrule
Example A & 24 & 0.320 +- 0.015\\
Example B & 36 & 0.415 +- 0.021\\
\bottomrule
\end{tabularx}
\par\smallskip\footnotesize\raggedright
All values are synthetic. Replace them with verified results, relevant
units and the actual uncertainty measure.
\end{table}

\section{Discussion}\label{sec:discussion}
Explain the benefits supported by the evaluation, limitations, appropriate
uses and maintenance plan. Distinguish current functionality from future
intentions. Refer to supporting Supplementary Figure~S1, Supplementary
Table~S1 or Supplementary Equation~(S1) from the separate supplement
where relevant.

\section*{Data Availability}
Replace with the working application URL, source-code and benchmark-data
access routes, versions and applicable licences. Include a permanent
archived code DOI when required at revision or acceptance. Confirm
maintenance and access commitments before stating them.

\section*{Supplementary Data}
Describe the supplied supplementary files and replace this instruction
with the journal's required availability statement when applicable.

\section*{Author Contributions}
Replace with the contributions agreed by every author, using appropriate
CRediT roles where applicable.

\section*{Acknowledgements}
Replace with confirmed acknowledgements, or remove if not applicable.

\section*{Funding}
Replace with confirmed funder names, grant identifiers and the source of
open-access funding where applicable.

\section*{Conflict of Interest}
Replace with the authors' confirmed competing-interest disclosure.

% Supply verified DOI and PMC links, or PubMed links when PMC is unavailable,
% in the corresponding BibTeX fields. Do not invent identifiers.
\bibliographystyle{prism-nar}
\bibliography{references}
\end{document}
